\begin{chapterbox}
    \chapter{Der Giftzwerg und das Drachenherz (2023)}
    \label{Der Giftzwerg und das Drachenherz (2023)}
    \az{71 und 72}

    \begin{center}
        Fortsetzung von \hypref{Der Giftzwerg und die Sphäre (2020)} und \hypref{Quruns Erwachen (2021)}
        
        einst versteckt unter dem Code 79082
    \end{center}
    
    Kjall, der diebische Sammler wertvoller Artefakte aus allen möglichen Zeiten, nutzt eine Konfrontation zwischen Varkur, Nomion und Shan. Er erringt Varkurs Zauberstab für sein geheimes Museum, lässt dabei aber auch spontan einen blutroten Edelstein mitgehen – und bereut diese Tat prompt, als \textit{Etwas} mit ihm Kontakt aufnimmt.
\end{chapterbox}

\az{Jahr 71}

Vorsichtig schlich sich Kjall an die Felsspalte heran und spähte in die Tiefe. Er erinnerte sich noch gut an den Vorfall, auch wenn es für ihn selbst schon einige Zeit her war. In der heutigen Nacht würde das gesamte Graue Gebirge gehörig durchgerüttelt werden. Den damaligen Gesprächen der Zwerge nach, die das Beben miterlebt hatten, würde ein gewaltiger Teil des verlassenen unterirdischen Zwergenreichs nahe der Korn-Schlucht einstürzen. Die Tiefminen würden besonders viele Flammen spucken, Kjall selbst würde sich nur knapp hinter einem Roteisenschild ducken und retten können. Einige seiner Kameraden würden nicht so glücklich sein.

Auch viele der höheren Schildzwerge waren erzürnt gewesen über die heutige Zerstörung uralter unterirdischer Bauwerke, auch wenn selbige bereits seit Jahrzehnten verlassen waren. Kjall sah die Lage ähnlich. Und ebenso vertraute er den ehemaligen Handwerkern der Schildzwerge, jahrhundertelang stabil bleibende Gänge zu bauen. Sie waren verlassen gewesen, weil das Volk der Schildzwerge leider schon lange nicht mehr so zahlreich war wie damals und weder alle Gänge vor finsteren Kreaturen verteidigen konnte noch wollte – was würde man mit all dem Extraplatz auch anfangen – aber sie waren definitiv nicht verlassen gewesen, weil Einsturzgefahr bestanden hätte.

Nein, Kjall stimmte den alten Rauschebärten zu, die felsenfest davon überzeugt waren, dass da Fremdeinwirkung im Spiel gewesen war. Fremdeinwirkung hieß, dass jemand Mächtiges sich eingemischt hatte und diese Gänge hatte einstürzen lassen. Und jemand Mächtiges im Grauen Gebirge bedeutete nun, wo Kjall viele Jahre älter war, dass er sich eventuell ein weiteres mächtiges Souvenir für sein geheimes Museum gönnen könnte. Schließlich war er nun ausgestattet mit der Sphäre, einem mächtigen Artefakt, welches es ermöglichte, Zeitportale in Vergangenes und in Zukünftiges zu öffnen,

So war Kjall in die Vergangenheit gehopst und schlich nun vorsichtig der Korn-Schlucht entlang, in der Hoffnung, den Urheber des Einsturzes ausfindig machen zu können. Und dem war so gewesen. Ein altbekanntes Schaudern hatte ihm schnell zu verstehen gegeben, dass ein mächtiges Wesen präsent war. Jemand, den Kjall kannte. Jemand, der sich hoffentlich nicht allzu sehr um Kjall scheren würde, da sie ihn noch gar nicht kannte. Kjalls alte, schattige Freundin, die einst als Teil von Kjalls Schar mit ihm in die Vergangenheit reisen würde, war offenbar heute in der Gegend unterwegs. Kjall bewegte sich vorsichtig in die Richtung, aus der dieses unwohle Gefühl stammte, welches Shan stets wie ein Schleier umgab. Dese Form von schleichender Kälte, die ihn an eine mondlose Nacht und an ein schattiges Tal erinnerte. Sie war unverwechselbar. Bald schon beobachtete Kjall, wie eine Masse aus festem Schatten in eine Felsritze ins Innere des uralten Zwergenreichs schlüpfte.

Kjall konnte seinen Augen kaum glauben, als er durch die kleine Felsspalte in einen tiefen Raum blickte und erkannte, wer dort unten an einem Blutsteinaltar stand, einen mächtigen roten Kristall in die Höhe hob, und verzweifelt mit Kjalls schattiger Freundin rangelte. Ein gewisser mächtiger Dunkler Magier, dessen Stab sich nur äußerst gut in Kjalls Kollektion machen würde. Von diesem roten Kristall ganz zu schweigen. Kjall rieb sich die schwieligen Hände.

Seit seinen ersten, unbeholfenen Zeitreisen hatte Kjall sich ein rigoroseres Protokoll zusammengestellt, wie er in einer solch lebensgefährlichen Situation vorgehen würde:

1. Wenn er sich selbst nicht aus der Zukunft auftauchen sah, würde er in die Vergangenheit reisen und irgendetwas Verrücktes, Aufsehenerregendes tun, das sein vergangenes Ich definitiv hätte wahrnehmen müssen (z.B. eine Steinflöte zücken und einen wilden Flötentanz aufzuführen beginnen).

2. Wenn er sich selbst aus der Zukunft auftauchen sah und seinem zukünftigen Ich etwas Unschönes geschah, so würde Kjall dessen Vorgehen genau beobachten, an diesen Moment zurückreisen und dann irgendetwas Aufsehenerregendes tun, das sein zukünftiges Ich definitiv nicht getan hatte (z.B. mitten in den Kampf zwischen seiner schattigen Freundin und dem Dunklen Magier springen und dabei laut die Melodie von Steinsang aus voller Kehle grölen).

3. Wenn er sich selbst aus der Zukunft auftauchen sah und sein zukünftiges Ich erfolgreich die lebensbedrohliche Gefahr umging und die zu erringenden Gegenstände errang, so würde Kjall dessen Vorgehen genau beobachten und dieses dann erfolgreich kopieren.

Aufgrund der Art und Weise, wie Zeitreisen mit der Sphäre von Cavern funktionierte, konnte Kjall ja nur in dieselbe Zeitlinie reisen, aus der er kam. Folglich konnte bloß dritte Fall eintreffen. Denn damit die Zeitreise-Logik nicht gebrochen wurde, musste Kjall sich wie sein zukünftiges Ich verhalten. Und Kjall hatte sich fest vorgenommen, die von seinem zukünftigen Ich beobachteten Handlungen nur zu kopieren, wenn der zukünftige Kjall glücklich mit dem Ergebnis war. So konnte er quasi das Schicksal erpressen: Ich verhalte mich nur so, wie du willst, wenn das Ganze so herauskommt, wie ich es will.

Ohne dass er je einen Finger dafür krümmen musste, wurde ihm quasi vom Schicksal selbst sein eigener bester Plan zur Ergatterung seltener Objekte aus Gefahrensituationen auf dem Silbertablett serviert. Was für ein Geschenk die Macht der Zeitreise doch war. Was für ein Narr er gewesen war, als er sich bei seiner ersten Zeitreise über die Unerschütterbarkeit der Zeitlinie geärgert hatte! Kjall kicherte. Damals war ihm dieses Korsett wie eine Limitation vorgekommen, das Xolls Leben vor ihm verwehrte, heute hingegen wie ein Geschenk. Solange er die Sphäre und genügend Treibstoff besaß, konnte ihm nichts Böses passieren.

Kjall hatte das Glückslos des Schicksals gezogen und lebte in der unglaublich unwahrscheinlichen Zeitlinie, in der er an die Sphäre gelangt war. Und die sich deshalb ganz nach seinen Wünschen entwickelte.

Sich strikt an diesen Plan haltend, blickte Kjall in die Höhle hinein, wo sich seine schattige Freundin und der Dunkle Magier gehörig um sich schlugen. Zaubersprüche, dunkler Rauch und Fetzen Fleisch gewordenen Schattens wirbelten umher. Kjall blickte aber nicht auf das Getümmel, sondern suchte nach einer Spur seines zukünftigen Selbsts. Tatsächlich erkannte Kjall kurz darauf durch den Felsspalt, wie sich tief unter ihm in der Höhle ein blaues Zeitportal öffnete und ein zukünftiger Kjall hinaustrat. Kjalls schattige Freundin und der mächtige Dunkle Magier waren zu sehr in ihren Kampf verwickelt, als dass sie sich groß um den Neuankömmling scherten. Vermutlich war er ihnen nicht einmal aufgefallen. Gerade schleuderte ein Schattententakel den roten Kristall aus der Hand des Magiers in eine Ecke des Raumes, wo bereits der Zauberstab des Magiers gelandet war. Kjall grinste und notierte sich den Zeitpunkt auf seinem Chronometer.

Sein zukünftiges Ich massierte seine Schläfen und orientierte sich. Es war nie angenehm, durch das Zeitportal zu reisen, aber immerhin schaffte Kjall es inzwischen, nicht mehr ohnmächtig zu werden bei kurzen Reisen durch den Zeit-Limbus zwischen den Zeitportalen, diesem endlosen Raum aus gleißendem blauem Licht, in dem es keine Schwerkraft zum Orientieren und keine Luft zum Atmen gab.

Dann zwinkerte der zukünftige Kjall dem jetzigen weit über ihm zu und marschierte schnurstracks auf den roten Kristall und den Zauberstab zu. Er ergriff beide, sprang zurück in das Zeitportal, aus dem er gekommen war, und das Portal verpuffte in blauen Wirbeln.

Mission erfüllt.

Kjall wandte sich grinsend ab und suchte in seinem Kopf schon nach dem kürzesten bekannten Weg zur Höhle, in der er sein zukünftiges Ich gerade hatte auftauchen sehen. Durch die Zeit mochte er reisen können, wie es ihm beliebte, doch den Ort seines Auftauchens musste er weiterhin manuell aufsuchen.

Er war frohen Mutes.

Auch dieses Vorhaben von ihm würde von Erfolg gekrönt sein.

Das Schicksal wollte es so.\bigskip




\az{Jahr 72}


Der zukünftige Kjall stolperte aus einem Zeitportal in seine eigene Zeit, fernab der kollabierten Gänge, und hielt prompt inne, um stolz seine beiden soeben errungenen Schätze zu betrachten.

Der schwarze Zauberstab des Dunklen Magiers war ein Meisterstück, auch wenn Kjall ihn selbst nicht verwenden konnte. Doch dieser rote Kristall war es, der Kjalls Interesse besonders geweckt hatte. Er war zwar nicht so groß wie das mächtige Drachenauge, das Fürst Kram – Kjall spuckte aus – in einigen Monaten in den Tiefen der Tiefminen finden würde. Beachtlich war es dennoch, insbesondere aufgrund der blutroten Farbe und des tiefschwarzen Schemens, der unter seiner Oberfläche umherzuwandern schien. Und jemand hatte zwar mehr schlecht als recht uralte Zwergenrunen in seine Oberfläche geritzt. Etwas Magisches hatte dieser Stein an sich, dessen war sich Kjall sicher.

Die glatte Oberfläche kribbelte unter seinen Fingern, fast, als würden diese ertauben. Rasch steckte Kjall den Edelstein in seinen Mantel.

Dann machte er sich auf den Weg in sein geheimes Museum tief in den Tiefminen. Die engen Gänge waren erfüllt von emsiger Bewegung. Dutzende Kjalls rasten umher, auf das Museum zu oder vom Museum weg, mit errungenen Schätzen in ihren Händen oder Plänen zur nächsten Errungenschaft in ihren Köpfen. Kjall hatte sich den heutigen Tag in dieser Zeit ausgesucht, um sein Museum in einem Schlag auf den neusten Stand zu füllen. Er fühlte sich ganz hibbelig, wie jedes Mal, wenn er seine Sammlung betrat. Rund um ihn herum eilten andere Kjalls, kreuzten seinen Weg oder nickten ihm freundlich zu, allesamt ähnlich grinsend wie er. Jeder einzelne von ihnen war er bereits gewesen oder würde er bald sein. Er fühlte sich wie in einem triumphalen Fiebertraum.

Dann war er da, bei seinem geheimen Museum. Er senkte seinen Kopf ehrfürchtig vor der großen Hammeraxt von Xoll Hammeraxt, seinem ehrenwerten Vater. Diese Waffe hatte über Kjalls Werkstatt gehangen, nun jedoch auch in seinem Museum. Sie hatte ihren Platz verdient neben Magischen Waffen, Mächtigen Schilden, Gaben aller Völker der bekannten Welt und dem einen oder anderen verschollenen Gegenstand.

Sechs Tiefminen-Golems, je zwei Versionen jeder seiner Kreationen, eilten im Museum umher, nahmen Schätze entgegen und deponierten sie sorgfältig auf Sockeln. Einmal mehr verfluchte Kjall Fürst Hallwort dafür, diese Wunderwerke der Mechanik, angetrieben von uralten Golemkernen, geschmiedet aus dem reisten Roteisenerz, in der Vergangenheit einfach eingeschmolzen zu haben. Dieser sonst so neugierige Dödel von einem Fürsten hatte das Roteisen wohl für eigene Zwecke nutzen wollen, statt diese Kunstwerke zu studieren. Elend! Doch immerhin hatte Kjall seine Golems jetzt wieder. Und nicht nur sie. Weitere metallene Helfer, solche, die für ihn vor langer Zeit in den Baum der Lieder eingebrochen waren und ihm erstmals die ersten Blätter des uralten Berichts zur Sphäre verschafft hatten, wuselten am Boden umher, beschrifteten die erjagten Schätze akribisch und signalisierten den verschiedenen Kjalls, wohin sie zu gehen hatten. Es war das reinste Chaos. Kjall würde keine Freude daran haben, dies alles zu organisieren, aber diese Aufgabe schob er fröhlich seinem zukünftigen Selbst zu. So wortwörtlich hatte er noch nie prokrastinieren können.

Gerade sah Kjall, wie einer seiner kleinen metallenen Helfer eine von Säure zerfressene Konstruktion auf einem Sockel deponierte. die in etwa die Form eines Arms hatte. Eine Prothese, die Kjall höchstpersönlich angefertigt hatte. Er knurrte grinsend beim Gedanken daran, sie erschaffen zu haben.

Angefangen hatte all dies mit der alten Graella aus den Tiefminen. Sie hatte ihre beiden Beine bei einem Feuerstoß aufgrund einer Grubengasexplosion verloren. Die Heiler hatten sie nur mit Stumpen zurückgelassen. Die Arme konnte so nicht mehr arbeiten, kein Teil der Kumpels mehr sein. Da hatte Kjall sich ein Herz gegriffen. Er kannte sich mit Tiefminen-Golems aus, mit mechanischen Armen und Beinen, und mit eigenständigen Läufern. Es war ein leichtes gewesen, der alten Graella metallene Beine zurückzugeben. Lange hatte er daran gearbeitet und sich an Graellas Reaktion erfreut. Gelegentlich musste sie nachölen, aber ihre neuen Füße waren sogar besser als ihre alten.

Doch dabei war es nicht geblieben. Als nächstes verlor der kleine Halbun einen Finger unter einem Hammer und fragte bei Kjall nach. Dann beschwerte sich ein Zwerg darüber, dass Graella schneller rannte als er selbst, und wollte besondere Zwergenstiefel von Kjall, die ihn gleich schnell machen würden. Dachten sie etwa, dass Kjall sich nur damit auseinandersetzte? Er war immer noch ein Tiefminen-Arbeiter! Er war ein Förderer von Bodenschätzen, ein Erforscher des Wissens, ein Sammler von Trophäen, kein Heiler für die Bedürftigen! Und die eigentlichen Heiler Caverns wollten nichts von seinen unorthodoxen Methoden hören, die verstanden sich bloß auf Fleisch und Blut. Nicht, dass dies Kjall gestört hätte, er wollte auch nicht irgendwelchen vorlauten Knirpsen seine Berufsgeheimnisse verticken. Er wollte einfach nur, dass er in Ruhe gelassen wurde. Die Altvorderen waren schließlich auch ohne moderne Prothesen ausgekommen.

Und es hatte nicht mit Zwergen geendet. Kjall spuckte aus. „Dies ist eindeutig kein Produkt aus zwergischen Werkstätten, dafür hat es zu viel Technik und zu wenig Magie in sich“, hatte diese hochnäsige Priesterin Gända vom Baum der Lieder abschätzig gesagt, als sie Graellas mechanische Beinkonstruktion inspiziert hatte – nur um einige Monate später dann dennoch diesen alten Wolfskrieger zu Kjall zu schicken. Wilselm hieß der Wolfskrieger, und einen ganzen Arm hatte er im Kampf gegen den Schwarzen Herold verloren. Und nun hatte er sich einen ganzen neuen Arm von Kjall gewünscht! Dabei hatte ja noch einen anderen, und er wollte ihn auch gar nicht, um zu kämpfen, sondern, weil es ihm praktischer im Alltag wäre. Was für ein Weichling. Und ein Mensch war es auch noch, einer aus Brandurs Schar, dieses Landräubers, der für Xolls Tod verantwortlich war.

Kjall kicherte bei der Erinnerung daran. Ja, er hatte dem Wolfskrieger Wilselm eine Konstruktion gebastelt. Nur nicht genau so, wie dieser wollte. Und schon bald hatte dieser sich nicht mehr danach erfreut.

Sein metallener Helfer lenkte Kjall zu einem freien Wandstück, in das ein Set aus Haken eingelassen war. Ehrfürchtig platzierte Kjall den Zauberstab des Dunklen Magiers darauf und trat einige Schritte zurück. Lange konnte er den Schatz jedoch nicht mehr bewundern, ehe er von seinem metallenen Helfer gleich wieder aus dem Raum gescheucht wurde, um Platz für weitere Tiefminen-Golems und Kjalls zu machen.

Kjall grinste, als er im Weggehehen seinen Blick über seine neusten Errungenschaften schweifen ließ: Eine andorische Flöte, die Brandur selbst gespielt hatte. Ein hadrischer Kompass, der den Seekrieger Ruuf damals zur mythischen Insel Danwar geleitet hatte. Boords Hammer, mit dem er die große Statue in der Halle der vier Schilde aus dem Stein gehauen hatte. Die Kapuze des Hral, die ihren Träger beinahe mit dem Schatten verschmelzen ließ. Ein Takuri-Spiegel aus der fernen Metropole Agarb.

So viele Möglichkeiten, was Kjall sich als nächstes beschaffen konnte.

Das Hauptaugenmerk des geheimen Museums waren natürlich die Sockel der vier mächtigen Schilde aus uralter Zeit. Drei hatte Kjall bereits erbeutet, und sie standen auf Hochglanz poliert nebeneinander auf ihren angestammten Sockeln: Der Sternenschild, der Bruderschild und der Silberschild. Nur der Feuerschild fehlte noch, stattdessen war auf dessen Sockel eine Skizze des Schwertmeisters Harthalt abgebildet, wie er den markanten dreieckigen Schild führte. Ob dieser Lackaffe überhaupt gewusst hatte, was für einen Schatz er da getragen hatte?

Dieser Schild würde Kjalls Sammlung vervollständigen, sein kleines Museum perfektionieren, seinem Lebenswerk die Krone aufsetzen. Sobald Kjall genug von seinen kleineren heutigen Errungenschaften hatte, würde er den Feuerschild aus der Vergangenheit holen und ins Museum bringen, aber erst am morgigen Tag. Diesen Moment der Perfektion und Erfüllung wollte er sich nicht durch ein Gewusel aus Kopien seiner Selbst, seiner Golems und seine metallenen Helfer stören lassen, und erst recht nicht wollte er sich diesen Moment, auf den er so lange hingearbeitet hatte, spoilern. Nein, den Feuerschild würde er allein holen, andächtig, nachdem er die Lust am ganzen Gewusel hier verloren hatte. Er war schon gespannt, wie viele Magische Waffen er bis dahin erringen konnte. Von Varlion dem Flammenschwert hatte er bereits eine Vermutung, wann er es erringen konnte, und Carlion das Kälteschwert stand sogar bereits in seiner Scheide in der hinteren rechten Ecke des Museums. Die Fürstenkrone der vier Schilde fehlte ihm ebenso wie die Rietgraskrone, doch die Krone der Nordmeere lag sicher verwahrt in einer staubfreien Vitrine hinter ihm. Mehr aus Spitzfindigkeit hatte er sogar den legendären goldenen Eintopftopf der Mutter des Zwergenfürsten Kram errungen, den er in einer Ecke als Spucknapf verwendete. Krams Mutter, wie hieß sie nun schon wieder? Murna? Oder doch eher Bairen? Egal, es war nicht wichtig.

Kjall selbst verstand nicht ganz, was hier alles abging, was für Artefakte hier alle noch in seinem geheimen Museum eintreffen und von emsigen metallenen Helfern einsortiert würden. Er war jedoch überzeugt davon, dass alles nach Plan lief.

Gedankenverloren zog er den roten Kristall hervor, den er in der Vergangenheit errungen hatte. Seine Fingerspitzen kribbelten. Dunkle Schemen pulsierten unter der Oberfläche. Was war dieses Objekt nur? Wo sollte er es hinstellen, wie sollte er es beschriften?

Direkt neben Kjall öffnete sich ein blaues Portal. Heraus stolperte eine weitere verschwitzte, doch grinsende Version seiner Selbst, die in ihrer Hand eine eiserne Kugel schwenkte, welche grünlich dampfte. Eine Kugel von Kreatoks Fallengas, direkt aus der großen Schmiede des Wunderzwergs höchstpersönlich! Leider wusste Kjall bereits, dass dieses Artefakt (oder zumindest sein Inhalt) bald schon wieder sein Museum verlassen mussten. Aber da ragte auch ein Fangzahn von Nehal höchstpersönlich aus der Tasche seines anderen Selbsts, und dieser würde ein schöner dauerhafter Bestandteil seiner Kollektion werden. Hach, die glorreichen Urzeiten des Zwergenreichs vor dem Unterirdischen Krieg. Vielleicht könnte Kjall dort Ferien machen gehen, wenn er hier fertig war. Er nickte seinem anderen Selbst anerkennend zu, welches zurückzwinkerte.

Aus einer dunklen Ecke von Kjalls Museum löste sich eine kleine Gestalt, die Kjall gerade mal bis zum Knie reichte. Ihr Körper war eine wabernde silberne Masse, die sich weigerte, scharfe Konturen anzunehmen. Nichtsdestotrotz glaubte Kjall, Arme und Beine ausmachen zu können, als dieses mysteriöse Wesen mit raschen Schritten auf das sich schließende Portal in die tiefe Vergangenheit zustürzte und darin verschwand. Es vergingen nur Augenblicke, bevor das Portal sich mit einem leisen Plopp wieder schloss.

Kjall blickte sich argwöhnisch seinem anderen Selbst nach, um zu sehen, was dieses davon hielt, doch war jenes bereits weitergelaufen und hatte sich in die Menge der Kjalls gemischt. Und länger konnte sich Kjall nicht mit diesem eigenartigen Vorfall befassen, denn in diesem Augenblick regte sich etwas tief in ihm. Seine Hand kribbelte und juckte. Kjall blickte argwöhnisch auf den roten Kristall, den er darin hielt, und befahl seinen Fingern, sich zu lösen und das Ding fallen zu lassen.

Seine Finger gehorchten ihm nicht mehr.

Kjall versuchte, sich zu rühren, um Hilfe zu schreien, irgendetwas zu tun, außer still und erstarrt im Raume zu stehen und regelmäßig ein- und auszuatmen. Es gelang ihm nicht. Da war eine gewisse Kälte, die durch seinen Arm pulsierte und bis in seinen Kopf vordrang. Kjall erschauderte. Da war Etwas, das sich gegen seinen Geist schmiegte. Gedanken, Erinnerungen und Gefühle wallten in ihm auf, die ihm unbekannt, fremd, falsch erschienen. Eine Präsenz war da, ein Dunkles Etwas, das ihn sondierte. Die Berührung dieser fremden Entität war schmerzhaft. Sie wirkte gleichzeitig zaghaft, als wäre das Wesen vorsichtig, und rabiat, als wäre das Wesen ungeübt in dem, was es hier tat. Und unzweifelhaft gingen all diese seltsamen Empfindungen vom roten Edelstein in Kjalls Händen aus.

Warum half ihm niemand? Warum schlug ihm kein zukünftiger Kjall den Kristall aus seinen Händen? Warum hatte ihn keiner gewarnt?

Eine tiefe Stimme scholl durch Kjalls Kopf, lauter als jede Stimme, die er in seinem Leben je gehört hatte. „WAS?! WIE KANN ICH ... WAS ... WER BIST DU?“

Kjalls ganzer Körper zitterte angespannt, doch abgesehen davon vermochte er, keinen Finger zu rühren. Die dumpfe Stimme schrie weiter: „AH, MEIN KRISTALLENES GEFÄNGNIS ... VARKUR GELANG ES BEREITS, MEINE GITTERSTÄBE ZU LOCKERN. EIN TEIL MEINER SELBST VERMAG ES, HERAUSZUFLIESSEN, UND DA BIST DU, UM MEINE GEDANKEN AUFZUFANGEN. WELCH WUNDER! WER BIST DU? EIN ZUKÜNFTIGER? EIN ZEITREISENDER? DIESE SPHÄRE ... DIESE MACHT!“

„Eine Macht, die dir nicht zusteht. Wer oder was bist du?“, dachte Kjall so laut, wie er konnte. Das Wesen beachtete ihn nicht, sprach jedoch immerhin erheblich leiser weiter.

„Du hast mich gerettet! Du hast mich vor Nomion bewahrt! Du hast mich Varkur und Shan entrissen! Du hast mir eine neue Zukunft ermöglicht. Du bist ...“ Der Druck auf Kjalls Schädel nahm zu. „... du bist Kjall, Sohn des Xoll. Und du hast einen Blick auf diese Welt, auf diese Zeit, wie ich selten einen wahrnahm. Wir können gemeinsam Großes erreichen, wir beide.“

Kjall ignorierte das Gefasel des Wesens. Er interessierte sich nicht dafür, wer dieser Nomion war, und er verfluchte Varkur und Shan gerade dafür, dass sie seine Aufmerksamkeit auf sich gezogen und ihn in diese Situation gebracht hatten. Kjalls Arm zuckte zur Seite, jedoch ohne den Kristall fallen zu lassen. Er hatte sich ein klein wehren können, entgegen dieser Beherrschung! Nun war er wieder erstarrt, doch wenn er sich nur ein klein wenig mehr anstrengte ...

Das Böse in ihm kicherte. „Doch zunächst einmal müssen wir uns an einen sicheren Ort begeben, wo ich diese neue Fähigkeit trainieren kann. Die Erstreckung meines Geistes auf Körper, die mein kristallenes Gefängnis berühren, das wird mir sehr zu Diensten sein ... und dass ich von allen Personen das Glück hatte, auf dich zu treffen, einen Zeitreisenden! Varkur und ich haben lange über diese Möglichkeiten gesprochen, damals in unserer gemeinsamen Zeit in Nordgard. Varkurs Mentor Koraph war äußerst fasziniert vom Studium der Zeit, auch wenn er diese nie so sehr wie Kirr beherrschte. Varkur berichtete von Hadrischen Stundengläsern, die die Zeit verlangsamen konnten, und davon, wie er und sein Mentor Koraph nächtelang darüber getratscht hatten, was es denn bedeuten würde, wenn man die Zeit nicht nur komplett anhalten, sondern vielleicht sogar so sehr verlangsamen könnte, dass sie rückwärts abliefe. Sie kamen zum eindeutigen Schluss, dass Zeitportale, wie sie sich beispielsweise in den Tiefen des Horun zu manifestieren scheinen, einen nur in dieselbe Zeitlinie befördern können, aus der man stammt. Und dass es theoretisch keinem bekannten Gesetz der Natur widerspräche, einen Abstecher ins Vergangene oder Zukünftige zu machen ... doch diese Sphäre, dieses Wunderwerk ... ich hätte es mir nie auszumalen gewagt, dass es so einfach wäre!“

Wie von selbst bewegten sich Kjalls Arme, verschoben Hebel und Schalter auf seiner Sphäre und öffneten ein rotierendes Zeitportal in eine einige Jahre zurückliegende Vergangenheit. Und auch wenn Kjall sich so stark wie möglich dagegen wehrte, konnte er nicht verhindern, dass sein Körper sich hineinwarf.





\az{Jahr 48}

Kjall wurde nicht ohnmächtig. Sein Körper plumpste durch das andere Zeitportal hinaus und schlug hart auf felsigem Boden der Vergangenheit auf. Das Portal wirbelte noch kurz vor sich hin, ehe es sich zusammenzog und mit einem leisen Plopp verschwand.

Dann war alles still.

Totenstill.

Noch immer spürte Kjall diese fremde Macht, die seinen Geist beherrschte. Wie ein dunkler Schemen am Rande seines Gesichtsfelds, der versuchte, selbiges zu überdecken und seinen Geist schlafen zu schicken. Doch noch schien dieser fremde Geist, überwältigt von der Zeitreise durch den Limbus, nicht die Macht dazu zu haben, ihn zu unterdrücken.

Kjall blickte sich um. Eine äußerst unangenehme Erfahrung, nicht bestimmten zu können, wohin seine Augen rollten. Dies war natürlich derselbe Raum, der sein geheimes Museum werden würde, doch befanden sich hier in der Vergangenheit noch keine errungenen Schätze darin. Noch nicht einmal das zertrümmerte Hadrische Stundenglas, welches Kjall nach der ersten Schlacht um die Rietburg ergattern würde. Stattdessen lagen auf den Sockeln erste aus Roteisen gefertigte Werkzeuge in allen erdenklichen Formen. Werke, auf die Kjall durchaus mit Freude zurückblicke – von ihrer Herstellung stammten viele der Brandnarben, die er mit Stolz trug. Dennoch waren diese Werke nichts im Vergleich zu den Tiefminen-Golems, die er später erschaffen würde.

„Welch ... widerliche Erfahrung, diese Reise durch den Zeit-Limbus“, sprach die finstere Macht mit seiner Kehle, jedes Wort verzerrt, als wäre es noch nicht an seine Stimmbänder gewöhnt. Kjalls Mund spuckte aus. Wie von selbst setzten sich Kjalls Beine in Bewegung, stolperten los durch die Tiefen der Tiefminen, auf Pfaden, die dieses Wesen aus seinen Erinnerungen gelesen haben musste, denn es verirrte sich nicht einmal auf feuerstoßgefährdeten Pfaden und ging nicht einmal einen Umweg.

Kjall hoffte darauf, dass ihm jemand entgegenkommen und ihn retten würde, irgendein Kumpel aus den Tiefminen, Radan, Drak, ja, selbst einer Begegnung mit Kram würde er erleichtert entgegensehen. Doch die Gänge blieben leer.

An einem kleinen Höhlenausgang irgendwo in den nördlichen Ausläufern des Grauen Gebirges kamen sie ins Freie. Ein mondloser Sternenhimmel beschien hohe Gipfel und vereinzelte Bäume. Die Finsternis der Nacht lag über den Gipfeln und Hügeln. Kalte Luft drang in seine Lungen und versorgte seinen vom Dauerlauf geschundenen Körper mit frischer Energie.

Kjall hatte inzwischen aufgegeben, mit dieser fremden Macht verhandeln zu wollen, die seinen Körper immer selbstsicherer dirigierte. Angst keimte in ihm auf, als er zu verstehen versuchte, warum kein zukünftiges Selbst ihn retten kam. War er von nun an verdammt, ein Sklave dieser fremden Macht zu sein? Waren alle zukünftigen Kjalls, die er in seinem geheimen Museum gesehen hatte, von diesem Wesen gesteuert worden?

„So“, sprach das Böse heiser in Kjalls eigener Stimme. „Du hast dich bislang beinahe brav an meine Anweisungen gehalten. Doch sind wir auch einige Male fast gestolpert, weil du dich mir zu widersetzen versuchtest. Wenn ich deinen Körper kontrolliere, dann doch bestimmt auch das weiche Ding, das in deinem Schädel herumschwimmt und dich schlafen lassen kann? Lassen wir uns doch ...“

Etwas pochte in seinem Kopf. Furcht wallte in Kjall auf, als eine unheimlich starke Müdigkeit aus dem Nichts durch seinen Körper schoss und seine Augen zufielen. Wenn dieses Ding seinen Geist einschlafen lassen und seinen Körper wie den eines Schlafwandlers fernsteuern könnte, wäre es erst recht um seinen Widerstand geschehen! Wobei, wenn er sich ohnehin nicht gegen sie wehren konnte, war ein Traum oder traumloser Schlaf womöglich Jahren der Agonie vorzuziehen.

Doch, ehe er eine Entscheidung fällen oder sich dem Mantel der Ungewissheit und Ohnmacht hingeben konnte, wurden Kjall und sein Tormentor durch ein fremdes, fernes Singen abgelenkt.

„Still“, knurrte das Böse, während es langsam zur Quelle dieser Geräusche schlich. Kjall, welcher ohnehin nicht hätte schreien können, horchte bloß. Das war nicht nur eine Stimme, das war ein ganzer dissonanter Chor, der in einer ihm fremden Sprache vor sich hin dudelte.

Dann lief Kjalls Körper um einen Felsen herum und sie beide erblickten in der Ferne die Urheber des gruseligen Gesangs.

Schwarz ragte der Dunkle Tempel in den grauen Himmel auf. Der riesige Turm war ein vielstöckiger Bau, in Kjalls eigener Zeit nur noch eine Ruine, von den rachsüchtigen Helden von Andor bis auf die Grundmauern eingerissen, um sein ätzendes Geisterfeuer einzudämmen.

Kjall wusste, dass düstere Drachenkultisten und geheimnisvolle Greifenanbeter der Grehon gelegentlich gemeinsam dunkle Messen in diesem geschichtsträchtigen Gemäuer abgehalten hatten. Genau wie jetzt. Eine Ansammlung an Gestalten in Dunklen Kapuzen hatte sich auf der Turmspitze versammelt, unterhalb eines gewaltigen ... war das ein Skelett? Befand sich ein riesiges Skelett auf der Turmspitze? Was auch immer diese Kultisten hier taten, Kjall hätte gerne so wenig wie möglich damit zu tun. Hoffentlich sah die fremde Entität dies ebenfalls so. Jeder konnte doch verspüren, dass unter dem Dunklen Tempel eine finstere Macht lauerte.

Doch das Böse in seinem Körper stolperte langsam näher, auf den überwachsenen Eingang des Dunklen Tempels hinzu, als wäre es magisch angezogen davon.

Und dann, als es nur noch einige Schritte davon entfernt war, geschah etwas Unerwartetes. Der gruselige Gesang der Kultisten wandelte sich in ekstatisches Geschreie. Ein dumpfes Dröhnen drang aus dem felsigen Boden. Hitze brandete auf. Und schwarz-silbernes Geisterfeuer brach an unregelmäßigen Stellen aus dem Boden rund um den Dunklen Tempel herum hervor und stieb als dunkel schimmernde Stichflammen in den Nachthimmel.

Kjall wusste nur von einem Artefakt, das schwarz-silbernes Feuer hervorrief, in welchem selbst Drachenknochen verkohlten. Nur die wenigsten kannten seine Geschichte, doch er hatte sie von seinem Vater Xoll schon oft erzählt bekommen. Dies war ... unerwartet. Doch erfreulich. Vielleicht musste Kjall sich gar nicht mit diesem Lackaffen Harthalt streiten, um den Feuerschild zu erringen. Voller Staunen und Ehrfurcht betrachtete Kjall das tödliche Feuerspiel vor ihm.

Und selbst die dunkle Macht, die seinen Körper kontrollierte, schien für einen Augenblick von einem ehrfürchtigen Gefühl der Bestimmung befallen zu sein. Seltsame Eindrücke schwammen durch Kjalls Geist, Bilder und Eindrücke, die er nicht einordnen konnte. Kristallsplitter sanken in einen See voller umherwirbelnder Farben, wie er sie noch nie zuvor gesehen hatte. Blutrote Flammen schmolzen wabernde Wolken von einem sternenübersäten Himmel. Ein gewaltiger Schmiedehammer presste einen krummen Turm in ein golden glühendes Feld.

„Es ... es ist wunderschön“, krächzte Kjalls Stimme. Für den flüchtigsten Augenblick verlor das fremde Ding den Halt über ihn. Auf einmal, wider allen Erwartens, gelang es Kjall, seinen Arm zum Zucken anzuregen. Er zwang seine Finger, sich zu öffnen und den roten Kristall loszulassen.

Klirrend fiel der Edelstein zu Boden. Kjall stolperte zitternd zurück, tastete sich fahrig ab, horchte in sich hinein. Keine Spur mehr des Bösen. Was auch immer in diesem Kristall gesteckt hatte, es hatte ihn nicht mehr in seiner Gewalt.

Das schwarz-silberne Feuer erlosch. Der Dunkle Tempel lag einen Herzschlag lang ruhig da, dann ertönte Jubeln und Klatschen von seiner Spitze, als die Kultisten einen nächsten Singsang anstimmten.

Kurz überlegte Kjall, sich dem roten Kristall wieder zu nähern. Dann jedoch schüttelte er seinen Kopf, ließ den Kristall Kristall sein und raste davon, bis er an einem stillen Örtchen ein Zeitportal in die Zukunft öffnen konnte. Was auch immer dies für eine dunkle Entität gewesen war, er hoffte, dass sie in der Vergangenheit ein baldiges Ende finden konnte. Ärger anrichten konnte dieses Ding in dieser Zeitlinie nicht mehr, das wäre ihn doch sicherlich schon aufgefallen.

Nein, er selbst hatte sich wichtigeren Dingen zuzuwenden.

Der Feuerschild wartete auf ihn.
